\documentclass[12pt,a4paper]{article}
\usepackage[utf-8]{inputenc}
\usepackage[french]{babel}
\usepackage{geometry}
\usepackage{fancyhdr}
\usepackage{xcolor}
\usepackage{listings}
\usepackage{amsmath}
\usepackage{amssymb}

% Configuration des marges
\geometry{top=2cm, bottom=2cm, left=2cm, right=2cm}

% Configuration des en-têtes et pieds de page
\pagestyle{fancy}
\fancyhf{}
\rhead{\textbf{Aide-mémoire C}}
\lhead{2025-2026}
\cfoot{\thepage}

% Configuration du code C
\lstset{
    language=C,
    basicstyle=\ttfamily\small,
    keywordstyle=\color{blue},
    commentstyle=\color{gray},
    stringstyle=\color{red},
    breaklines=true,
    showstringspaces=false,
    frame=single,
    numbers=left,
    numberstyle=\tiny,
    stepnumber=1,
    tabsize=4
}

\title{\textbf{Aide-mémoire pour l'examen C}}
\author{Nelen Nathan}
\date{\today}

\begin{document}

\maketitle
\tableofcontents
\newpage

\section{Bases essentielles de C}

\subsection{Entrées et sorties}

\begin{lstlisting}
#include <stdio.h>

// Affichage
printf("Message : %d\n", valeur);
printf("Texte : %s\n", chaine);

// Lecture
scanf("%d", &entier);
scanf("%f", &reel);
scanf("%s", tableau_char);
\end{lstlisting}

\textbf{Formats courants :}
\begin{itemize}
    \item \texttt{\%d} : entier
    \item \texttt{\%f} : nombre décimal
    \item \texttt{\%s} : chaîne de caractères
    \item \texttt{\%p} : adresse (pointeur)
\end{itemize}

\subsection{Déclaration de variables}

\begin{lstlisting}
int nombre = 5;
float reel = 3.14;
char caractere = 'A';
char chaine[] = "Bonjour";
\end{lstlisting}

\section{Pointeurs}

\subsection{Concept fondamental}

Un pointeur est une \textbf{variable qui stocke l'adresse mémoire} d'une autre variable.

\begin{lstlisting}
int x = 5;
int *ptr = &x;    // ptr contient l'adresse de x

printf("%d\n", x);      // affiche 5
printf("%d\n", *ptr);   // affiche 5 (valeur pointée)
printf("%p\n", ptr);    // affiche l'adresse de x
\end{lstlisting}

\subsection{Règles essentielles}

\begin{itemize}
    \item \textbf{\texttt{\&x}} : adresse de la variable \texttt{x}
    \item \textbf{\texttt{*ptr}} : valeur stockée à l'adresse contenue dans \texttt{ptr}
    \item \textbf{\texttt{int *p}} : déclare un pointeur vers un entier
\end{itemize}

\subsection{Modifier une valeur via un pointeur}

\begin{lstlisting}
int x = 10;
int *p = &x;

*p = 20;  // Modifie x à travers p
printf("%d\n", x);  // affiche 20
\end{lstlisting}

\section{Fonctions}

\subsection{Syntaxe générale}

\begin{lstlisting}
// Déclaration (prototype)
type nom_fonction(parametres);

// Définition
type nom_fonction(parametres) {
    // code
    return valeur;
}

// Appel dans main
int main() {
    nom_fonction(arguments);
    return 0;
}
\end{lstlisting}

\subsection{Passage par valeur vs pointeur}

\begin{lstlisting}
// Passage par valeur (ne modifie pas l'original)
void modifier_valeur(int x) {
    x = 100;
}

// Passage par pointeur (modifie l'original)
void modifier_pointeur(int *p) {
    *p = 100;
}

// Utilisation
int main() {
    int a = 5;
    
    modifier_valeur(a);      // a reste 5
    modifier_pointeur(&a);   // a devient 100
    
    return 0;
}
\end{lstlisting}

\section{Fonctions récurrentes dans vos exercices}

\subsection{Somme et Produit d'un tableau}

Fichier : \texttt{Labo\_Tableau-Pointeur/Exercice\_Tableau-Pointeur.c}

\begin{lstlisting}
void sommeEtProduit(int *tab, int taille, 
                    int *somme, int *produit) {
    *somme = 0;
    *produit = 1;
    
    for (int i = 0; i < taille; i++) {
        *somme += tab[i];
        *produit *= tab[i];
    }
}

// Appel dans main
int main() {
    int tab[] = {1, 2, 3, 4, 5};
    int somme, produit;
    
    sommeEtProduit(tab, 5, &somme, &produit);
    printf("Somme = %d\n", somme);
    printf("Produit = %d\n", produit);
    
    return 0;
}
\end{lstlisting}

\subsection{Échange de deux valeurs}

Fichier : \texttt{Pointeur/1-Pointeur/Exercice\_Pointeur.c}

\begin{lstlisting}
void echange(int *a, int *b) {
    int temp = *a;
    *a = *b;
    *b = temp;
}

// Utilisation
int main() {
    int x = 5, y = 10;
    
    echange(&x, &y);
    printf("x = %d, y = %d\n", x, y);  // x = 10, y = 5
    
    return 0;
}
\end{lstlisting}

\subsection{Logarithme népérien}

Fichier : \texttt{Labo\_Fonction/main.c}

\begin{lstlisting}
float LogNeperien(float x, float n) {
    float lnX = 0, exposant, lnProvisoire;
    
    exposant = 1;
    for (int i = 1; i <= n; i++) {
        lnProvisoire = pow((x-1)/(x+1), exposant);
        lnProvisoire *= 1/exposant;
        lnX += lnProvisoire;
        exposant += 2;
    }
    
    return lnX * 2;
}
\end{lstlisting}

\section{Tableaux}

\subsection{Déclaration et initialisation}

\begin{lstlisting}
// Tableau vide (taille fixe)
int tab1[10];

// Tableau initialisé
int tab2[] = {1, 2, 3, 4, 5};
int tab3[5] = {1, 2, 3, 4, 5};

// Tableau multidimensionnel
int matrice[3][3] = {
    {1, 2, 3},
    {4, 5, 6},
    {7, 8, 9}
};

// Chaîne de caractères (tableau de char)
char chaine[] = "Bonjour";
\end{lstlisting}

\subsection{Accès et modification}

\begin{lstlisting}
int tab[] = {10, 20, 30};

tab[0] = 10;  // Accès par indice
printf("%d\n", tab[1]);  // affiche 20

int *ptr = tab;  // Pointeur sur le tableau
printf("%d\n", *ptr);     // affiche 10
printf("%d\n", *(ptr+1)); // affiche 20
\end{lstlisting}

\subsection{Passage de tableau à une fonction}

\begin{lstlisting}
void afficher_tableau(int tab[], int taille) {
    for (int i = 0; i < taille; i++) {
        printf("%d ", tab[i]);
    }
    printf("\n");
}

// Appel
int main() {
    int tab[] = {1, 2, 3, 4, 5};
    afficher_tableau(tab, 5);
    return 0;
}
\end{lstlisting}

\section{Structures}

\subsection{Déclaration et utilisation}

\begin{lstlisting}
// Déclaration
struct Personne {
    int age;
    char nom[30];
    float taille;
};

// Utilisation dans main
int main() {
    struct Personne p1;
    p1.age = 25;
    strcpy(p1.nom, "Alice");
    p1.taille = 1.70;
    
    printf("Age : %d\n", p1.age);
    printf("Nom : %s\n", p1.nom);
    
    return 0;
}
\end{lstlisting}

\subsection{Tableau de structures}

\begin{lstlisting}
struct Personne personnes[3];

personnes[0].age = 25;
strcpy(personnes[0].nom, "Alice");
\end{lstlisting}

\section{Boucles et conditionnelles}

\subsection{Boucles}

\begin{lstlisting}
// Boucle for
for (int i = 0; i < 10; i++) {
    printf("%d\n", i);
}

// Boucle while
int i = 0;
while (i < 10) {
    printf("%d\n", i);
    i++;
}

// Boucle do-while
int j = 0;
do {
    printf("%d\n", j);
    j++;
} while (j < 10);
\end{lstlisting}

\subsection{Conditions}

\begin{lstlisting}
if (age >= 18) {
    printf("Majeur\n");
} else if (age >= 13) {
    printf("Adolescent\n");
} else {
    printf("Enfant\n");
}

// Switch
switch (choix) {
    case 1:
        printf("Option 1\n");
        break;
    case 2:
        printf("Option 2\n");
        break;
    default:
        printf("Option invalide\n");
}
\end{lstlisting}

\section{Opérateurs importants}

\begin{tabular}{|c|c|c|}
\hline
\textbf{Opérateur} & \textbf{Description} & \textbf{Exemple} \\
\hline
\texttt{+} & Addition & \texttt{a + b} \\
\texttt{-} & Soustraction & \texttt{a - b} \\
\texttt{*} & Multiplication & \texttt{a * b} \\
\texttt{/} & Division & \texttt{a / b} \\
\texttt{\%} & Modulo & \texttt{a \% b} \\
\hline
\texttt{==} & Égal & \texttt{a == b} \\
\texttt{!=} & Différent & \texttt{a != b} \\
\texttt{<} & Inférieur & \texttt{a < b} \\
\texttt{>} & Supérieur & \texttt{a > b} \\
\hline
\texttt{\&\&} & ET logique & \texttt{a \&\& b} \\
\texttt{||} & OU logique & \texttt{a || b} \\
\texttt{!} & NON logique & \texttt{!a} \\
\hline
\end{tabular}

\section{Conseils d'examen}

\begin{enumerate}
    \item \textbf{Déclarez les fonctions avant \texttt{main()}}
    \item \textbf{N'oubliez pas le \texttt{\&}} avec \texttt{scanf}
    \item \textbf{Attention aux pointeurs} : \texttt{\&x} vs \texttt{*p}
    \item \textbf{Initialisez les variables} avant utilisation
    \item \textbf{Nommez les fonctions} clairement et logiquement
    \item \textbf{Testez avec des valeurs simples} pour déboguer
    \item \textbf{Compilez souvent} pour détecter les erreurs tôt
\end{enumerate}

\section{Checklist finale}

\begin{itemize}
    \item[\checkmark] Inclusion des headers (\texttt{\#include <stdio.h>}, etc.)
    \item[\checkmark] Prototype des fonctions avant \texttt{main}
    \item[\checkmark] Pointeurs correctement déclarés avec \texttt{*}
    \item[\checkmark] Paramètres passés correctement (\texttt{\&} ou adresse)
    \item[\checkmark] Pas de variables non initialisées
    \item[\checkmark] Boucles et conditions correctement écrites
    \item[\checkmark] \texttt{return 0;} à la fin de \texttt{main}
\end{itemize}

\end{document}
